%% ==========================================================================
%%  8  Computational Evidence
%%
%%  Includes an explicit statement of why the randomised evidence is weaker
%%  than its raw counts suggest.  Do not delete that subsection: the sampler
%%  bias is the reason the search has found nothing, and building adversarial
%%  generators is the highest-value experiment outstanding.
%% ==========================================================================
\section{Computational Evidence}
\label{sec:experiments}

All searches operate directly on the graph-theoretic form of
Conjecture~\ref{conj:main}: enumerate allocations, build the envy graph, detect
positive-weight cycles and compute $\pathw{\alloc}$ by a Bellman--Ford
iteration on maximum-weight walks, and record
$\min_{\alloc} \max_i \pathw{\alloc}(i)$.

\begin{center}
\begin{tabular}{@{}p{0.46\textwidth}lp{0.30\textwidth}@{}}
\toprule
Experiment & Script & Result \\
\midrule
Exhaustive: $n = m = 3$, all $9{,}880$ instances up to agent symmetry
  & \texttt{search.py} & worst-case minimum per-agent subsidy $= 1$ \\
Randomised: $n \in \set{3,4,5}$, $m \in \set{4,5,6}$, ${\approx}\,3.5\times10^4$
instances
  & \texttt{fast.py} & no counterexample \\
Binary additive construction of Theorem~\ref{thm:binadd}, $2\times10^5$
instances, $n \le 5$, $m \le 7$
  & \texttt{binadd.py} & maximum per-agent subsidy ever observed $= 1$ \\
The obstruction of Theorem~\ref{thm:obstruction}
  & \texttt{deadend.py} & all three insertions require subsidy $2$; full
    instance solvable with $0$ \\
Does \emph{some} utilitarian-optimal allocation attain ${\le}\,1$?
  & \texttt{msw.py} & yes, in all but one of ${\approx}\,2.6\times10^4$
    instances \\
Does \emph{every} utilitarian-optimal allocation attain ${\le}\,1$?
  & \texttt{msw.py} & no --- fails in ${\approx}\,2.5\%$ of instances at
    $n=3, m=4$, rising to ${\approx}\,5\%$ at $m = 5$ \\
The single exception
  & \texttt{mswcex.py} & refutes Conjecture~\ref{conj:msw}
    (Proposition~\ref{prop:msw-false}) \\
\midrule
Peel schedule and trap-state refusal on the witness of
Theorem~\ref{thm:obstruction}
  & \texttt{peel.py} & six legal peels reach the zero-subsidy allocation; the
    trap state has $\pathw{}(1) = 2$ \\
Full peel state space of that witness, $7^3 = 343$ states, both move types
  & \texttt{peel3.py} & 175 legal, 166 good, \textbf{9 dead ends}, root good,
    all six reachable terminals are perfect matchings \\
\bottomrule
\end{tabular}
\end{center}

The exhaustive row is the only unconditional evidence for
Conjecture~\ref{conj:main}. The dichotomous functions on $m$ items number
$2, 6, 38, 990$ for $m = 1,2,3,4$. Agents are interchangeable, so an instance is
a multiset of three of them, and at $m = 3$ there are
$\binom{38 + 3 - 1}{3} = \binom{40}{3} = 9{,}880$ of these --- feasible to
enumerate, and already out of reach at $m = 4$.

The last three rows are the ones that changed the shape of the problem. The
utilitarian-optimal set almost always contains a witness, but not always, and
one exception is enough: Proposition~\ref{prop:msw-false} is built from the
single instance in which it did not.

% --------------------------------------------------------------------------
\subsection{Why the randomised evidence is weak}
\label{sec:samplerbias}

The count of ${\approx}\,3.5 \times 10^4$ instances without a counterexample
overstates the support for Conjecture~\ref{conj:main}, and the reason is
structural rather than a matter of sample size.

The generator builds a dichotomous function by visiting subsets in order of
increasing size and, at each subset $S$, computing the interval
$[\ell, h]$ of values consistent with what has already been fixed --- namely
$\ell = \max_{g \in S} \cost(S \setminus \set{g})$ and
$h = \min_{g \in S} \cost(S \setminus \set{g}) + 1$ --- then choosing an
endpoint by an unbiased coin whenever $\ell \ne h$. Independent fair coins put
almost all of the mass on functions that wander in the middle of the feasible
region. Structured extremal functions, which take the same endpoint at every
subset, are exponentially unlikely: on $m$ items there are $2^m - 1$ nonempty
subsets and a function that requires a prescribed endpoint at each of them is
sampled with probability $2^{-(2^m-1)}$, which is $1/128$ at $m = 3$ and about
$3 \times 10^{-5}$ at $m = 4$, per agent.

This matters because the two instances that carry the results of this paper are
of exactly that kind. The obstruction of Theorem~\ref{thm:obstruction} uses
$\cost_1(S) = \max(0, \abs{S}-1)$ alongside $\cost_2 = \cost_3 = \abs{S}$, all
three of which are endpoint-constant; drawing that triple at $m = 3$ has
probability $2^{-21} \approx 5 \times 10^{-7}$, so a run of $3.5 \times 10^4$
samples would be expected to produce it about once in every sixty runs. Both
instances were found by hand or by a targeted search, not by the sampler. The
sampler has, in effect, never looked in the region where the interesting
behaviour lives.

\begin{remark}[The experiment worth running]
\label{rem:adversarial}
If Conjecture~\ref{conj:main} is false, the counterexample is in the corner the
current generator cannot reach, and no amount of additional uniform sampling
will find it. An adversarial generator --- biased towards endpoint-constant
functions, towards supermodular costs, towards agents whose free chores overlap
heavily, and towards instances in which the utilitarian optimum is unique --- is
the single highest-value experiment outstanding against the \emph{conjecture}.
It should be built before much further effort is spent on proof attempts,
because it is the only cheap way to learn whether there is a theorem to prove.
\end{remark}

% --------------------------------------------------------------------------
\subsection{The peel-frame search}
\label{sec:peelsearch}

The highest-value experiment against the \emph{approach} of
\Cref{sec:approach3} is different, and it is cheap. Reuse the
\texttt{gen.py}/\texttt{search.py} harness over the full $n = 3$, $m = 3$
dichotomous family --- all $9{,}880$ instances up to agent symmetry --- and for
each instance run the fixed-point computation of \texttt{peel3.py}: is the root
good, and how many dead ends are there? An instance with a bad root refutes
Conjecture~\ref{conj:h1prime} and kills the peel approach.

It would \emph{not}, by Remark~\ref{rem:converse}, refute
Conjecture~\ref{conj:main}: a good terminal state can in principle exist while
every schedule reaching it passes through an illegal intermediate state. The two
searches therefore test different things and both are worth running --- the
allocation-space search above tests the conjecture, and the peel-frame search
tests whether this route to it survives.

Two follow-ups, in order. First, test the hypothesis of
Remark~\ref{rem:balance} across that family: is every dead end a state already
committed to an unbalanced terminal? If so, formulate the balance rule precisely
and check whether greedy-under-that-rule reaches a terminal without
backtracking. Second, try candidate rules in the order (a) balance-first, peel
so that a balanced terminal remains reachable; (b) spectator-free peels first by
Lemma~\ref{lem:specfree}, ties broken by balance; (c) $\argmax_i \pathw{}(i)$
with a balance tie-break --- noting that (c) alone already fails on the witness.

% --------------------------------------------------------------------------
\subsection{Reproducing the results}
\label{sec:repro}

\begin{center}
\begin{tabular}{@{}ll@{}}
\toprule
Script & Purpose \\
\midrule
\texttt{gen.py}        & enumerate all dichotomous functions on $m$ items \\
\texttt{search.py}     & exhaustive search, arguments \texttt{m n} \\
\texttt{fast.py}       & randomised counterexample hunt, \texttt{m n T seed} \\
\texttt{msw.py}        & utilitarian-optimal test, \texttt{m n T seed} \\
\texttt{tiebreak.py}   & selection rules of \Cref{sec:tiebreaks} \\
\texttt{localsearch.py}& local minima of single-transfer search \\
\texttt{deadend.py}    & Theorem~\ref{thm:obstruction} \\
\texttt{binadd.py}     & Theorem~\ref{thm:binadd} \\
\texttt{mswcex.py}     & Proposition~\ref{prop:msw-false} \\
\texttt{peel.py}       & the schedule and trap-state refusal of
                         \Cref{sec:dissolve} \\
\texttt{peel3.py}      & the state-space table of Proposition~\ref{prop:deadends} \\
\bottomrule
\end{tabular}
\end{center}

The instance of Proposition~\ref{prop:msw-false} was produced by
\texttt{msw.py 4 3 3000 11} at trial $2460$ and is re-derived independently by
\texttt{mswcex.py}, which re-checks that all three cost functions are negative
dichotomous, that the utilitarian optimum is unique, and that a zero-subsidy
allocation exists outside it.
